\documentclass{article}
\usepackage{mathtools} 
\usepackage{fontspec}
\usepackage[UTF8]{ctex}
\usepackage{amsthm}
\usepackage{mdframed}
\usepackage{xcolor}
\usepackage{amssymb}
\usepackage{amsmath}


% 定义新的带灰色背景的说明环境 zremark
\newmdtheoremenv[
  backgroundcolor=gray!10,
  % 边框与背景一致，边框线会消失
  linecolor=gray!10
]{zremark}{说明}

% 通用矩阵命令: \flexmatrix{矩阵名}{元素符号}{行数}{列数}
\newcommand{\flexmatrix}[4]{
  \[
  #1 = \begin{pmatrix}
    #2_{11}     & #2_{12}     & \cdots & #2_{1#4}   \\
    #2_{21}     & #2_{22}     & \cdots & #2_{2#4}   \\
    \vdots      & \vdots      & \ddots & \vdots     \\
    #2_{#31}    & #2_{#32}    & \cdots & #2_{#3#4}
  \end{pmatrix}
  \]
}

% 简化版命令（默认矩阵名为A，元素符号为a）: \quickmatrix{行数}{列数}
\newcommand{\quickmatrix}[2]{\flexmatrix{A}{a}{#1}{#2}}

\begin{document}
\title{注释 4.2}
\author{张志聪}
\maketitle
\begin{zremark}
  在区间上连续函数的值域也是一个区间。即值域也是连续的。
\end{zremark}

\textbf{证明：}

函数是常值函数，则值域是单点，也是区间。

设$f$在区间$I$上连续（不是常值函数），令
\begin{align*}
  V = f(I)
\end{align*}

任取$u, v \in V$且$u < v$，由函数定义可知，存在$a, b \in I$使得
\begin{align*}
  f(a) = u, f(b) = v
\end{align*}
接下来，我们要证明$[u, v] \subseteq V$。

因为$[a, b] \subseteq I$，所以$f$在$[a, b]$（或$[b, a]$）上连续，由介值定理可知
任意$u < y_0 < v$都有$x_0 \in (a, b)$（或$[b, a]$）使得
\begin{align*}
  y_0 = f(x_0)
\end{align*}
所以，任意$y \in [u, v]$都属于$V$，所以$[u, v] \subseteq V$。
由$u, v$的任意性可知，$V$是一个区间。

\begin{zremark}
  函数$f$在区间$I$上连续且可逆，那么$f$在$I$上严格单调。
\end{zremark}

\textbf{证明：}

对任意$x_1 < x_2$，由$f$是可逆函数可知，$f$是单射函数，
于是可得
\begin{align*}
  f(x_1) \neq f(x_2)
\end{align*}

假设$f$不是严格单调的，那么存在
$x_1 < x_2 < x_3 \in I$，有
\begin{align*}
  f(x_1) < f(x_2), f(x_2) > f(x_3)
\end{align*}
或
\begin{align*}
  f(x_1) > f(x_2), f(x_2) < f(x_3)
\end{align*}
（图像出现局部“转折”）。以下的证明按第一种情况处理，另一种情况证明类似。

由介值定理可知，取$f(x_1) < y_0 < f(x_2), f(x_3) < y_0 < f(x_2)$，
存在$x_0 \in (x_1, x_2), x_0' \in (x_2, x_3)$使得
\begin{align*}
  f(x_0) = y_0 \\
  f(x_0') = y_0
\end{align*}
这与函数单射矛盾，假设不成立，命题得证。


\begin{zremark}
  函数$f(x)$在区间$I$上连续，那么$|f(x)|$也在$I$上连续。
\end{zremark}

\textbf{证明：}

对任意$x_0 \in I$，由$f(x)$在$I$上连续可知，
对任意$\epsilon > 0$，存在$\delta > 0$，当$|x - x_0| < \delta$时，有
\begin{align*}
  |f(x) - f(x_0)| < \epsilon
\end{align*}
因为，我们有
\begin{align*}
  | |f(x)| - |f(x_0)| | \leq |f(x) - f(x_0)| < \epsilon
\end{align*}
所以$\lim\limits_{x \to x_0} |f(x)| = |f(x_0)|$，
由$x_0$的任意性可知，$|f(x)|$也在$I$上连续。

\begin{zremark}
  $f, g$在区间$I$上一致连续，则$f + g, f - g$也在$I$上一致连续，
  $f g, f / g$（$g(x) \neq 0$）不一定一致连续。
\end{zremark}

\textbf{证明：}

\begin{itemize}
  \item $f + g$

        由$f, g$在$I$上一致连续可知，
        对任意$\epsilon > 0$，存在$\delta_1 > 0$，使得任意$x^\prime, x^{\prime\prime} \in I$，
        只要$|x^\prime - x''| < \delta_1$，就有
        \begin{align*}
          |f(x^\prime) - f(x^{\prime\prime})| < \epsilon
        \end{align*}

        同理，存在$\delta_2 > 0$，使得任意$x^\prime, x^{\prime\prime} \in I$，
        只要$|x^\prime - x''| < \delta_2$，就有
        \begin{align*}
          |g(x^\prime) - g(x^{\prime\prime})| < \epsilon
        \end{align*}

        取$\delta = \min(\delta_1, \delta_2)$，对任意$x^\prime, x^{\prime\prime} \in I$，
        只要$|x' - x''| < \delta$，就有
        \begin{align*}
          |f(x') + g(x') - [f(x'') + g(x'')]|
           & = |f(x') - f(x'') + g(x') - g(x'')|      \\
           & \leq |f(x') - f(x'')| + |g(x') - g(x'')| \\
           & \leq \epsilon + \epsilon = 2\epsilon
        \end{align*}
        所以，$f + g$在区间$I$上一致连续。

  \item $f - g$

        易得$-g$在$I$上是一致连续的，然后可得$f + (-g)$在区间$I$上一致连续。

  \item $fg$

        举个反例，$f(x), g(x) = x$在$\mathbb{R}$上一致连续，但$(fg)(x) = x^2$在区间$\mathbb{R}$上不是一致连续的。
        $\S 2$习题13

  \item $f/g$

        举个反例，$f(x) = 1, g(x) = x$在$(0, +\infty)$上一致连续，但$(f/g)(x) = \frac{1}{x}$在区间$(0, +\infty)$上不是一致连续的。
        $\S 2$例9

\end{itemize}

\begin{zremark}
  函数$f$在$[a, b]$上一致连续，那么$f$在$(a, b]$上一致连续。
\end{zremark}

\textbf{证明：}

按照定义，结论是显然的。

\begin{zremark}
  设$f$在$[a, +\infty)$上连续，且$\lim\limits_{x \to +\infty} f(x)$存在，
  则$f$在$[a, +\infty)$上一致连续。
\end{zremark}

\textbf{证明：}

习题16


\end{document}